% Original predict-observe-explain lesson. Counts are ideal Born probabilities, not experimental data.
\documentclass[12pt,aspectratio=43]{beamer}
\usetheme{Boadilla}
\usecolortheme{orchid}
\usepackage{amsmath,braket,quantikz}
\setbeamertemplate{navigation symbols}{}
\title{A qubit is not a hidden bit}
\subtitle{Predict, observe, explain}
\author{Quantum information tutorial}
\date{Tutorial 07}
\begin{document}
\begin{frame}[plain]\titlepage\end{frame}
\section{Predict}
\begin{frame}[label=prediction]{What comes out of the second Hadamard?}
\centering
\begin{quantikz}
\lstick{$\ket{0}$} & \gate{H} & \gate{H} & \meter{}\end{quantikz}
\bigskip
\begin{enumerate}
\item Always 0.
\item Half 0, half 1.
\item Always 1.
\end{enumerate}
\note{Take a show of hands. Ask whether the first gate is being treated as a random classical flip.}
\end{frame}
\begin{frame}{Compute the amplitudes before the probabilities}
\[
 H=\frac{1}{\sqrt2}\begin{pmatrix}1&1\\1&-1\end{pmatrix},
 \quad H\ket{0}=\frac{\ket{0}+\ket{1}}{\sqrt2}.
\]
\pause
\[
 H^2=I,\qquad H^2\ket{0}=\ket{0}.
\]
\begin{exampleblock}{Ideal prediction}
Without an intermediate measurement, the final computational-basis outcome is 0 with probability one.
\end{exampleblock}
\end{frame}
\section{Intervene}
\begin{frame}{Now insert an intermediate measurement}
\begin{overlayarea}{\textwidth}{4.2cm}
\only<1>{Measure after the first $H$, then apply the second $H$. Predict the final distribution again.}
\only<2>{The first measurement yields $\ket{0}$ or $\ket{1}$, each with probability $1/2$. Applying $H$ to either gives equal computational-basis probabilities.}
\only<3>{The final distribution is now $P(0)=P(1)=1/2$. Measurement changed the state preparation for the remaining circuit.}
\end{overlayarea}
\end{frame}
\begin{frame}{A phase can matter without changing one measurement}
\[
 \ket{+}=\frac{\ket0+\ket1}{\sqrt2},\qquad
 \ket{-}=\frac{\ket0-\ket1}{\sqrt2}.
\]
\begin{columns}[T]
\begin{column}{.47\textwidth}\begin{block}{Measure now}
Both give 0 and 1 with equal probabilities in the computational basis.
\end{block}\end{column}
\begin{column}{.47\textwidth}\begin{block}{Apply $H$ first}
$H\ket{+}=\ket0$, while $H\ket{-}=\ket1$.
\end{block}\end{column}
\end{columns}
\end{frame}
\section{Explain}
\begin{frame}{State vectors carry more than a classical histogram}
\begin{itemize}[<+->]
\item Relative phase affects interference under a later operation.
\item Probabilities are squared moduli of amplitudes.
\item A measurement basis is part of the prediction.
\end{itemize}
\[
 \ket\psi=\alpha\ket0+\beta\ket1,\qquad |\alpha|^2+|\beta|^2=1.
\]
\end{frame}
\begin{frame}{Two qubits: prepare a Bell state}
\centering
\begin{quantikz}
\lstick{$\ket0$} & \gate{H} & \ctrl{1} & \qw\\
\lstick{$\ket0$} & \qw & \targ{} & \qw
\end{quantikz}
\[
 \ket{\Phi^+}=\frac{\ket{00}+\ket{11}}{\sqrt2}.
\]
\small Computational-basis results are correlated. This single basis of data alone does not establish entanglement against every alternative preparation.
\end{frame}
\begin{frame}{Exit ticket: distinguish three statements}
\begin{description}
\item[State] What mathematical object represents the preparation?
\item[Operation] Which transformation occurs before measurement?
\item[Observation] Which basis and outcome probabilities are being reported?
\end{description}
\bigskip
Return to the first circuit and explain why ``two random flips'' gives the wrong model.
\end{frame}
\begin{frame}{Continue the tutorial}
\small
\href{https://quantum.cloud.ibm.com/learning/en/courses/basics-of-quantum-information/single-systems/introduction}{IBM Quantum Learning: Basics of Quantum Information} provides background on state vectors, measurements and operations.
\medskip
The lesson sequence and circuits here are original; all probabilities assume ideal gates and measurements. There are no hardware execution results in this deck.
\end{frame}
\end{document}
